

%قبل از شروع مسئله ترجیحاً چیزی ننویسید. اما اگر خواستید بنویسید، در قالب لیستی (فقرات و شمارش) بنویسید.
%برای اضافه کردن نمره، بعدی از شروع‌مسئله، [x] را اضافه کنید و به جای x  نمره مورد نظر را قرار دهید. برای زیرمسئله‌ها هم به همین صورت است.
%\شروع‌مسئله[10]
\شروع‌مسئله []
%\متن‌سیاه{سوال عملی اول:}\\
%\متن‌سیاه{دراپ‌تست}\\

شما $n$ نقطه روی صفحه‌ای کشیده‌اید و بازی‌ای را با شنگدباو شروع می‌کنید. او هر سری به شما دو عدد $u,v$ می‌گوید و شما باید یال بین $u$ و $v$ را روی کاغذ بکشید. اگر گرافی که با یال‌های کنونی‌ای که کشیده‌اید دوبخشی نبود باید بگویید «بینگو» و بازی تمام می‌شود.

\متن‌سیاه{ورودی:}\\
در خط اول به شما دو عدد $n , m$ داده میشود در خط $i$ ام $i$ امین یالی که قرار است شما اضافه کنید قرار دارد
و به صورت دو عدد $u , v$ آمده است که به این معناست که شما باید یالی بین $u , v$  اضافه کنید

$$1 \le n, m \le 10^5$$

\متن‌سیاه{خروجی:}\\
تنها یک عدد $i$ را چاپ کنید که این به این معناست که پس از افزودن $i$ امین یال دیگر گراف دوبخشی نیست و شما باید بگویید «بینگو».

\متن‌سیاه{محدودیت‌ها:}\\
$1 \leq n, q \leq 10^{6}$\\
$0 \leq |s[i]| \leq 10^{9}$\\
$1 \leq a, b \leq n$\\

\متن‌سیاه{ورودی و خروجی نمونه:}\\

\begin{latin}
	\begin{tabularx}{\textwidth}{|X | X|}
		\hline
		\textbf{stdin} &
		\textbf{stdout}\\
		\hline
		4 6 & 5\\
		1 2 &\\
		3 4 &\\
		3 2 &\\
		1 4 & \\
		1 3&\\
		2 4&\\
		\hline
	\end{tabularx}
\end{latin}

\solution{
می‌دانیم از یک جایی به بعد دارای دور فرد شده و دوبخشی نمی‌شود. ما باید اولین جاییکه دور فرد ایجاد شده را پیدا کنیم. می‌توانیم جست‌وجوی دودویی روی این مکان بزنیم و به ازای هر اندیس مورد جست‌وجو تمامی یال‌ها تا قبل از آن اندیس را در نظر می‌گیریم. با استفاده از یک 
\lr{dfs}
می‌توان فهمید که آیا گراف دارای دور فرد هست یا نه. به اینصورت که رنگ هر کس را یا صفر یا یک قرار می‌دهیم و در 
\lr{dfs}
هر رأس را رنگ مخالف پدرش می‌گذاریم. می‌دانیم درخت 
\lr{dfs}
مجبور می‌کند گراف اینطوری رنگ شود پس اگر رنگ‌آمیزی ارائه شده توسط جست‌وجوی عمق اول ما مطلوب نبود یعنی یالی داشت که دو سرش رنگ برابر داشتند یعنی دوری فرد داریم. چک کردن چنین چیزی معادل چک کردن هر یال است در کل الگوریتم
$\mathcal{O}(nlog_2n)$
طول می‌کشد.

}
%}
% صورت اصلی مسئله را در این‌جا بنویسید.
%اگر مسئله چند قسمت داشت به این صورت عمل کنید:
%\شروع{زیرمسائل}
%\زیرمسئله[0] % مسئلۀ ۱آ
%\زیرمسئله % مسئلۀ ۱ب
%\زیرمسئله % مسئلۀ ۱پ
%\زیرمسئله % مسئلۀ ۱ت
%\پایان{زیرمسائل}



%%%%%%%%%%%%%%%%%%%%%%%%%%%%%%%%%%%%%%%%%%%%%%%%%%%%%
% برخی دستورات مفید برای تمرین‌ها. %
%%%%%%%%%%%%%%%%%%%%%%%%%%%%%%%%%%%%%%%%%%%%%%%%%%%%%

% فرمول‌های نمونه:
%$f(x) = \Theta(x)$
%$T(x, y) \leq \log(x) + 2^y + \binom{2n}{n}$
% {\tt code\_function}


% لیست موارد بدون شماره‌گذاری:
%\شروع{فقرات}
%    \فقره مورد اول 
%        \شروع{فقرات}
%            \فقره مورد اول داخلی 
%                \شروع{فقرات}
%                    \فقره اولین مورد خیلی داخلی 
%                    \فقره دومین مورد خیلی داخلی 
%                \پایان{فقرات}
%            \فقره مورد دوم داخلی 
%        \پایان{فقرات}
%    \فقره مورد دوم 
%\پایان{فقرات}

% لیست موارد شماره‌گذاری شده:
%\شروع{شمارش}
%    \فقره یک مورد 
%    \فقره یک مورد دیگر 
%    \فقره باز هم یک مورد دیگر
%\پایان{شمارش}

% برای نوشتن یک فرمول چند خطی که یک شماره بخوره. برای حالت‌های دیگر به سایت زیر مراجعه کنید:
% https://www.overleaf.com/learn/latex/Aligning_equations_with_amsmath
%\begin{align} 
%    \begin{split}
%        (x+y)^3 &= (x+y)^2(x+y) \\
%                &= (x^2+2xy+y^2)(x+y) \\
%                &= (x^3+2x^2y+xy^2) + (x^2y+2xy^2+y^3) \\
%                &= x^3+3x^2y+3xy^2+y^3
%    \end{split}                                 
%\end{align}

% نوشتن ماتریس و آرایه دوبعدی
%\begin{align}
%    A = 
%    \begin{bmatrix}
%        A_{11} & A_{21} \\
%        A_{21} & A_{22}
%    \end{bmatrix}
%\end{align}

% اضافه کردن تصویر با شرح. برای اضافه کردن تصویر بدون شرح، فقط دستور درج‌تصویر را استفاده کنید.
%\شروع{شکل}[h!]
%\تنظیم‌ازوسط
%  \درج‌تصویر[پهنا=0.5\پهنای‌متن]{img.jpg}
%\شرح{شرح تصویر}
%\برچسب{تصویر:۱}
%\پایان{شکل}
%\پایان{زیرمسائل}

%کد زیر را ببینید.
%%\lstinputlisting{source_filename.py}
%\begin{latin}
%\begin{lstlisting}
%# Import the modules
%import sys
%import random
%
%ans = True
%
%while ans:
%question = raw_input("Ask the magic 8 ball a question: (press enter to quit) ")
%
%answers = random.randint(1,8)
%
%if question == "":
%sys.exit()
%
%elif answers == 1:
%print "It is certain"
%
%elif answers == 2:
%print "Outlook good"
%
%elif answers == 3:
%print "You may rely on it"
%
%elif answers == 4:
%print "Ask again later"
%
%elif answers == 5:
%print "Concentrate and ask again"
%
%elif answers == 6:
%print "Reply hazy, try again"
%
%elif answers == 7:
%print "My reply is no"
%
%elif answers == 8:
%print "My sources say no"
%\end{lstlisting}
%\end{latin}

	
	